\pdfbookmark[0]{Anzeigen. Herschel: Satz von Cotes. 1814.}{anzeigen-herschel-1814}
\begin{center}
\rule{.85\textwidth}{.3pt}\par
Göttingische gelehrte Anzeigen.\quad 14. Februar 1814.\par
\rule{.85\textwidth}{.3pt}
\end{center}
\hspace*{\parindent}\textit{Philosophical Transactions of the Royal Society of London for the Year} 1813. IV und 304 Seiten. 26 Seiten \textit{Meteorological Journal} und 8 Seiten \textit{Index} im Quartformat.

\textit{Mathematische und astronomische Abhandlungen. --- Über eine bemerkenswerte Anwendung des Satzes von} \textsc{Cotes}, \textit{von} \textsc{J. F. W. Herschel} (Sohn des Astronomen). Es sei \(n\) eine beliebige ganze Zahl, \(n\omega=360^0\) und \(N\) ein beliebiger Winkel. Unter diesen Voraussetzungen gibt der Satz von \textsc{Cotes} durch einen einfachen Ausdruck das Produkt aller Radiusvektoren an, denen auf einem Kegelschnitt, astronomisch gesprochen, die wahren Anomalien \(N,\ N+\omega,\ N+2\omega,\ N+3\omega\mathpunct{.}\mathpunct{.}\mathpunct{.}\mathpunct{.}N+(n-1)\omega\) entsprechen. Diese und andere ähnliche Entwicklungen, die den Gegenstand des Aufsatzes bilden, bieten an sich keine besonderen Schwierigkeiten. Der Aufsatz ist wegen der Art der Behandlung dennoch mit Vergnügen zu lesen. Einige neuere mathematische Schriftsteller verwenden die Schreibweise \(\cos^2 A\) für das Quadrat von \(\cos A\), ganz gegen jede Analogie. Nach dieser müsste sie den Kosinus eines Bogens \(=\cos A\) bezeichnen. Was der Verfasser über diese Schreibweise sagt, findet unsere volle Zustimmung.

